\documentclass[ecta,nameyear,draft]{econsocart}
\usepackage{amsmath,amsfonts,amssymb,amsthm,booktabs,array,longtable}
\usepackage[T1]{fontenc}
\usepackage{mathpazo,microtype,needspace,xurl}
\usepackage{xr-hyper}
\RequirePackage[colorlinks,citecolor=blue,linkcolor=blue,urlcolor=blue]{hyperref}
\startlocaldefs
\newtheorem{theorem}{Theorem}
\newtheorem{proposition}[theorem]{Proposition}
\newtheorem{lemma}[theorem]{Lemma}
\newtheorem{corollary}[theorem]{Corollary}
\newcommand{\E}{\mathbb E}
\newcommand{\R}{\mathbb R}
\endlocaldefs
\setlength{\emergencystretch}{3em}
\makeatletter\def\copyright@text{Prepared for referee review}\makeatother

\setcounter{secnumdepth}{2}

\externaldocument[][nocite]{build/supp}[supp.pdf]
\externaldocument[][nocite]{build/appendix}[appendix.pdf]
\begin{document}
\begin{frontmatter}
\title{Neural Bellman Operators}
\runtitle{Neural Bellman Operators}
\begin{aug}
\author[id=au1,addressref={add1}]{\fnms{Qian}~\snm{QI}\ead[label=e77ECTA]{qiqian@pku.edu.cn}}
\address[id=add1]{Peking University}
\end{aug}
\begin{abstract}
This paper develops Neural Bellman Operators for constructive policy evaluation and feasible improvement. Witness-preserving continuations and qualification-optimal readouts connect attained approximation to the original Bellman optimum. Exact conditional integration and independently implemented whole-bin evaluation qualify the original-law costs of complete controllers. Joint marginal and paired intervals imply sharp controller-choice and deployment-allocation certificates, including installation costs and finite-cohort implementation. In the nonlinear-investment experiment, attained ReLU controllers have lower expected cost than matched Bernstein controllers at both acquisition resolutions. The joint intervals identify the cost ordering and conditional acquisition-fee regions. At an unresolved refinement-fee boundary, allocation reduces worst-case expected loss, subject to an explicit additional-setup threshold. These comparisons concern constructed continuations, not the optional neural query proposals whose adverse complete-work results are retained. Recursive, vector, predictable-qualification and witnessed-reuse results retain their assumptions, proofs and evidence.
\end{abstract}
\begin{keyword}
\kwd{Dynamic programming}\kwd{Neural approximation}\kwd{Policy certification}
\kwd{Information acquisition}\kwd{Computational economics}
\end{keyword}
\end{frontmatter}
\section{Introduction}\label{sec:introduction65}
A numerical continuation matters because it helps produce an economic decision. The error of a fitted value, the loss of its implemented policy and the resources required to verify that loss are different objects. A useful method must connect approximation and resource choices to a feasible decision under the original economic law.

This paper develops Neural Bellman Operators (NBOs) along that chain. The constructive implementation evaluates each continuation against its own fitted future and retains feasible action witnesses. Its exact affine--ReLU realization and primitive error budget do not require unrestricted nonconvex training to succeed. A complementary recursive implementation queries the original Bellman problem where action search requires information. A trained proposal can change query locations, but verified endpoint bounds and protected conventional refinement determine validity. These are distinct uses of neural approximation; success of a common certificate is not evidence that every neural component improves complete work.

The investment economy has state-dependent capacity, bilinear transitions and one or two independent continuous innovations. Primitive curvature bounds the transition's convexity defect and supplies regularity of the original Bellman value. Recursive certification can dispense with a stored state lattice, but not with the descendant cost of an original-value query. A strong-convexity upper chord can certify a rounded action before another endpoint query; the method pays for safeguarded refinement whenever the screen fails. A fixed-feature readout has a convex certificate-aware objective and an auditable optimization gap. Neither a neural continuation nor an arbitrary proposal is assumed convex.

The attained continuation provides a separate route to original-law cost evaluation. Its residual oscillations and implemented greedy gaps control both the observation range and original-optimum policy loss. The terminal correction allows an inexact fitted terminal value, and state-independent shifts cancel. For shallow affine--ReLU continuations, conditional expectations under the original one- and two-innovation laws are finite positive-part moment calculations. Explicit state-and-action covers qualify residuals without querying the optimal Bellman tree. Conditional on the actually fitted future, qualification-optimal readouts minimize that verified account; exact primal--dual checks certify the achieved gap rather than relying on solver status. Backward construction then returns the complete acquired controller that is evaluated.

Purchased original-value information can also be reused. A lower Bellman bound and a feasible upper witness are transported only after paying for state displacement, capacity repair and action rounding. Failed checks invoke the original verifier. The finite-cover fill bound states its retention and precision premises explicitly; the executed bounded cache receives only the guarantees its actual contract supports. Predictable qualifications and time-uniform cost intervals separately address when to stop purchasing information. Fixed-controller costs and conditional average costs of changing controllers are not interchanged, and failed numerical evaluations retain their qualified support and paid work.

The evidence distinguishes these constructions. The retained action-readout catalogue compares certificate-aware ReLU and equally sized Bernstein proposals, independent-context validation, sixty-four-block reuse and trained two-control extensions. The recursive-controller cost comparisons meet their declared equivalence margin without identifying strict superiority; adverse cold-start, fitting and reuse outcomes remain visible. Localized witnessed-reuse workloads count first fills and every descendant on a miss. They are not policy-cost observations or a long-horizon neural speed result.

The later original-law experiment evaluates the attained continuation actors themselves. ReLU and Bernstein constructions at each acquisition resolution use nine coefficients, common limits, primitive information and a matched acquisition contract. Their actual fitted futures are propagated backwards. Two estimators of the same complete-policy costs use one frozen simultaneous family: direct cash cost and exact last-innovation conditional integration. The resulting intervals identify lower ReLU-controller costs at both matched resolutions and order the five constructed controllers. This is a positive finite comparison between specified attained fits, not an ordering of all conventional or neural methods.

The economic decision uses those same costs. Sharp difference-constraint bounds retain all paid marginal and paired information when choosing a controller at declared resource prices. Deployment allocation extends the decision when neither acquisition precision is uniformly best after its fee. The controller is selected before state observation, so a share means a mixture of complete feasible policies rather than averaged actions or network coefficients. The new allocation theorem gives an attained linear-programming certificate, exact setup-cost thresholds and an implementable finite-cohort choice. It remains valid for data-dependent prices on the original coverage event. No sample count is relabeled as a deployment horizon, and installation, qualification and processor time are not assigned invented monetary values.

\paragraph{Relation to existing methods.}
Lipschitz extensions and distance transforms underlie the constructive envelope \citep{mcshane1934,felzenszwalb2012}. Approximate and safe policy iteration have established error analyses \citep{munos2003,pirotta2013,scherrer2015}, as does convex dynamic-programming approximation \citep{yang2020convex}. Learned warm starts and optimization-based control supply related numerical methods \citep{sakaue2023warm,sambharya2023,agrawal2020control}. Value-function control variates and telescoping identities are established evaluation tools \citep{munos2006variance,jiang2016}. We connect these ideas to a verified residual-oscillation account for an attained own-future continuation, continuous-law numerical enclosures and the implemented controller's cost.

Time-uniform concentration and off-policy confidence sequences have their own theory \citep{howard2020time,karampatziakis2021}. Our observations are direct original-law contrasts, and the contribution identifies the paid numerical qualification needed to use those concentration tools. Minimax-regret decisions have an established economic role \citep{manski2004}; shortest-path difference constraints are classical \citep{floyd1962}. Our deployment result uses the common attaining cost potentials of the paid interval system to obtain an exact allocation and installation account. It is conditional on qualified costs, not a new unconditional statistical-risk theorem.

\paragraph{Reading path.}
The main article develops construction, recursive and vector certification, qualified continuation fitting, predictable evaluation, actual controller costs and deployment choice. The Technical Supplement contains the core proofs. The Numerical and Development Appendix retains every earlier active section, the sharp pure-choice theorem, complete allocation proofs and all adverse numerical records. Historical companion editions retain controlled diffusions, recursive and endogenous preferences, temporal selves and games under their own hypotheses; the present expected-cost deployment identities do not silently extend to those different objectives.

\section{The economic operator and its policy account}\label{sec:core49}
Consider dates $t=0,\ldots,T$ and an invariant compact state domain $X\subset\R^d$. At date $t$, the action belongs to a nonempty compact set $A_t(x)$. The transition is $F_t(x,a,Z_{t+1})$, the current cost is $c_t(x,a)$, and the terminal cost is $g(x)$. Write $B_0=1$ and $B_t=\prod_{j<t}\beta_j$, with $\beta_j>0$. An implemented policy includes its observation rule, action selection and feasible numerical realization, not just its stored parameters.

For a continuation $v$, define
\begin{equation}\label{eq:operator49}
 Q_t(v;x,a)=c_t(x,a)+\beta_t\rho_t(v(F_t(x,a,Z))),\qquad
 \mathcal T_tv(x)=\inf_{a\in A_t(x)}Q_t(v;x,a).
\end{equation}
Throughout the deterministic comparison, the evaluations $\rho_t$ are monotone and cash invariant on a common domain containing all functions and constant translates used below. They are therefore nonexpansive in the supremum norm. Continuous objectives, effective compact feasible correspondences and measurable minimizers, or measurable arbitrarily near-minimizers with their error charged, are assumed. The stipulated conditional recursion gives Bellman comparison against all admissible adapted controls. In the experiments $\rho_t$ is conditional expectation, costs are bounded and innovations have their original continuous uniform law. The direct-cost inference in Section~\ref{sec:decisions49} is specific to expectation; it is not automatically an inference theorem for nonlinear recursive utility or games.

A Neural Bellman Operator is a construction and evaluation procedure for \eqref{eq:operator49}. It forms targets using its \emph{own} fitted $f_{t+1}$, returns a neural continuation and a feasible policy, and records an economic error allowance. It does not receive the optimal continuation as a training label. Arbitrary trained networks can be evaluated by the first result below. The subsequent constructive backend specifies how a network and its actor are obtained without assuming successful nonconvex optimization.

\subsection{Action-dependent errors and complete-policy loss}
Let $H(a)$ be a true current cost and $\widehat H(a)$ its continuation approximation. If $\widehat a$ is $\delta$-optimal for $\widehat H$ on the same feasible action set, then
\begin{equation}\label{eq:center49}
 H(\widehat a)-\inf_aH(a)\le\delta+
 \operatorname{osc}_{a}\{\widehat H(a)-H(a)\}.
\end{equation}
Indeed, compare $\widehat a$ with an arbitrarily near-optimal true action and cancel the approximate objectives. The infimum limit gives the result. Additive continuation levels cancel; action-dependent discrepancies do not. Under a change in future policies, the continuation must be transported to the new economic problem before this inequality licenses reuse. The full development edition retains the original quantitative transport, recursive-domain and economic-application results. No transport bound is presumed valid after an unaccounted change of constraints or preferences.

\begin{proposition}[One-sided policy sandwich]\label{prop:sandwich49}
Suppose, uniformly in state,
\[
 f_t-\mathcal T_tf_{t+1}\le u_t,\qquad
 \mathcal T_t^\pi f_{t+1}-f_t\le v_t,
 \qquad \ell_T\le f_T-g\le u_T.
\]
Then the actual feasible policy satisfies
\begin{equation}\label{eq:sandwich49}
 0\le J_0^\pi-V_0^*\le
 \sum_{t<T}B_t(u_t+v_t)+B_T(u_T-\ell_T).
\end{equation}
The bound is invariant to date-specific constant shifts of the fitted continuations, provided the associated residual accounts are shifted consistently.
\end{proposition}
\noindent The proof is in Supplement, Section~\ref{supp:proof-core49-1}.


A sharper valid residual account for the same continuation and the same deployed actor sharpens knowledge of the policy. It does not retrain that policy or change its realized cost. Conversely, subtracting two upper bounds in \eqref{eq:sandwich49} does not identify the difference in actual costs. These distinctions govern both the historical comparisons and the new experiments.

\section{Constructing and compiling a feasible neural policy}\label{sec:construction49}
Assume the primitive objective has, for feasible state--action pairs, state modulus $A_t^Q=C_{x,t}+\beta_tF_{x,t}L_{t+1}$ and action modulus $D_t=C_{a,t}+\beta_tF_{a,t}L_{t+1}$. Let
\[
 d_H(A_t(x),A_t(z))\le K_t^A\|x-z\|_1,
 \qquad L_t=A_t^Q+D_tK_t^A,\qquad L_T=L_g.
\]
The Hausdorff condition and the primitive moduli imply that $\mathcal T_tf_{t+1}$ is $L_t$-Lipschitz whenever $f_{t+1}$ is $L_{t+1}$-Lipschitz. All moduli and numerical integration procedures are supplied by the economic primitives, not fitted from successful outcomes.

At state nodes $x_{t,i}$ with covering radius $h_t$, form effective feasible $k_t$-nets. Evaluate their own-future objectives with certified error at most $e_t$. Minimize the computed labels and retain the corresponding \emph{original} action $a_{t,i}$. A Borel repair $\Gamma_{t,i}(x)$ is feasible at $x$ and satisfies
\begin{equation}\label{eq:repair49}
 \|\Gamma_{t,i}(x)-a_{t,i}\|_1\le K_t^A\|x-x_{t,i}\|_1+\zeta_t.
\end{equation}
The allowance $\zeta_t$ concerns displacement; feasibility itself may not fail. Define
\begin{equation}\label{eq:cone49}
 f_t(x)=\min_i\{y_{t,i}+L_t\|x-x_{t,i}\|_1\},\quad
 i_t(x)=\min\arg\min_i\{y_{t,i}+L_t\|x-x_{t,i}\|_1\}.
\end{equation}
Return $\pi_t(x)=\Gamma_{t,i_t(x)}(x)$. Terminal labels approximate $g$ with error $e_T$ and use the same cone construction with modulus $L_g$.

\begin{theorem}[Own-future construction with feasible witnesses]\label{thm:construct49}
The functions in \eqref{eq:cone49} have exact affine--ReLU realizations and Lipschitz moduli $L_t$ independent of label error. The returned policy is feasible and measurable. It satisfies
\begin{align}\label{eq:enclosures49}
 -e_t\le f_t-\mathcal T_tf_{t+1}&\le D_tk_t+e_t+2L_th_t,\\
 \mathcal T_t^\pi f_{t+1}-f_t&\le e_t+D_t\zeta_t,\notag
\end{align}
and hence
\begin{equation}\label{eq:constructionloss49}
 J_0^\pi-V_0^*\le
 \sum_{t<T}B_t(D_tk_t+2e_t+2L_th_t+D_t\zeta_t)
 +B_T(2e_T+2L_gh_T).
\end{equation}
With effective covers, repairs and certified queries at the prescribed accuracy, allocating each nonnegative term a share of a positive economic tolerance gives a terminating finite construction. A fixed numerical precision or a capped finite experimental ladder need not contain that allocation.
\end{theorem}
\noindent The proof is in Supplement, Section~\ref{supp:proof-core49-2}.


The Lipschitz-extension principle is classical \citep{mcshane1934}. A native min-plus envelope and its exact ReLU circuit, using the same labels, witnesses and ties, are \emph{one method with two encodings}. They return the same function and policy. Neither Theorem~\ref{thm:construct49} nor a faster representation proves a uniquely neural capability or a training-success probability for Adam or L-BFGS. The mathematical increment is the complete feasible, own-future construction and implementation account; the empirical question is its work and policy value relative to conventional generators.

\subsection{Original-owner compilation on a tensor cover}
On a tensor cover $\mathcal X=\prod_{j=1}^d\{x_{j,0},\ldots,x_{j,n_j}\}$, put $S=\prod_j(n_j+1)$. For arbitrary labels $y_i$, including inconsistent inexact labels, compute at each grid node $v$
\[
 (z_v,o_v)=\min_{\rm lex,i}(y_i+L\|v-x_i\|_1,i).
\]
The index $o_v$ is the original label owner, not the transformed node $v$.

\begin{theorem}[Continuous original-owner identity]\label{thm:compile49}
For a closed tensor cell $C$ with corners $\mathcal V(C)$, and $x\in C$,
\begin{align}\label{eq:compile49}
 f(x)&=\min_{v\in\mathcal V(C)}(z_v+L\|x-v\|_1),\\
 (f(x),i_*(x))&=\min_{\rm lex,v\in\mathcal V(C)}
 (y_{o_v}+L\|x-x_{o_v}\|_1,o_v).\notag
\end{align}
Coordinatewise forward and backward relaxations compute $(z,o)$ in $O(dS)$ additions and comparisons and $O(S)$ storage. After locating the cell, a query uses $O(d2^d)$ arithmetic operations. Exact rational comparisons preserve all ties, including $L=0$. Using the original owner's repaired action preserves the uncompiled policy, its cost and its certificate exactly.
\end{theorem}
\noindent The proof is in Supplement, Section~\ref{supp:proof-core49-3}.


The nodal algorithm is the classical distance transform \citep{felzenszwalb2012}; preserving the original action through continuous off-grid queries is the policy obligation. If there are $J_{t,i}$ node actions and $M_t$ integration representatives, write $Q_t=M_t\sum_iJ_{t,i}$. Literal dense cone evaluation costs $O(\sum_t dS_{t+1}Q_t)$ in representation work. Compilation replaces this by $O(\sum_t[dS_{t+1}+Q_td2^d])$, plus cell location, primitive evaluation, integration, feasible repair, certificate and output. Rational bit lengths and state boxes crossing many cells incur additional costs. There is no dimension-free complexity assertion.

\subsection{Acquired states and conventional actors}
A $b$-bit coordinate sensor reports a dyadic cell and its midpoint $\widehat x$. On the unit cube, $\|x-\widehat x\|_1\le r(b)=d2^{-b-1}$. For the investment capacity $\kappa(x)$ below, clip the original action to the smallest capacity throughout that cell and round down to an action quantum $\Delta_t$. This is feasible at every state in the cell. Relative to an exact true-state repair, its displacement is at most $2K_t^Ar(b)+\Delta_t$. Transporting the witness cone from $\widehat x$ to $x$ adds at most $2L_tr(b)$, plus any certified selector objective error $\nu_t$. Thus the implemented-policy account adds
\begin{equation}\label{eq:acquire49}
 \sum_{t<T}B_t\{(2L_t+2D_tK_t^A)r(b)+D_t\Delta_t+\nu_t\}.
\end{equation}
No continuity of the discrete owner selector is assumed. At uncertain numerical comparisons the simulator retains all possible actors, rather than choosing a favorable one.

For conventional FVI the continuation is multilinear on its realized tensor grid. Its actual modulus is bounded by the largest coordinate edge slope. A nearest-node actor uses a feasible node minimizer and true-state repair. Let $h_t$ be the largest half-cell $\ell^1$ radius, $L_t^{\rm gr}$ the Bellman graph modulus, and $e_t$ the node query allowance. Its optimal residual upper bound is $D_tk_t+e_t+L_t^{\rm gr}h_t$. Its selected-policy upper bound is $e_t+E_t$, where
\[
 E_t=\sup_x\{y_{i(x)}+L_t^{\rm gr}\|x-x_{i(x)}\|_1-f_t(x)\}.
\]
Partitioning cells by nearest-node boundaries makes the expression separately affine in each coordinate. Its maximum is therefore at a partition vertex; the implementation encloses all such vertices. Both uniform and curvature-refined FVI receive this sharper one-sided account, not an artificially doubled baseline. Their acquired-state allowance is derived separately in \eqref{eq:fvisensor49}.

\section{Original-optimum accuracy in the investment economy}\label{sec:accuracy62}
The exact witness theorem answers an implementation question conditional on a continuation. An economist also needs to know how close the complete implemented policy is to the original optimum. We now establish that connection in the investment economy itself. The argument exploits a property not used by the earlier cell-range calculation: despite the bilinear transition, the Bellman objective remains jointly convex. This permits a second-order interpolation account and transfers verified nodal decisions to an implementable policy over the entire state domain. It does not assume that a fitted neural continuation is convex.

\paragraph{Relation to the preceding two-control query.}
The earlier query experiment used the fixed shared capacity $1/8$. The complete dynamic extension below uses the original state-dependent capacity $c(x)$, with the retained exposure columns, action-cost interaction and second shock. Its nesting claim concerns the original scalar economic model; it is not an assertion that a full dynamic policy was already present in the earlier isolated query.

\subsection{Primitives and Bellman regularity}
Let $d\ge2$ be even, let indices be cyclic, and set $\beta=15/16$. Write
\begin{align*}
 S(x)&=\frac2d\sum_j(x_j-5/8)^2+
       \frac1{4d}\sum_j(x_j-x_{j+1})^2+
       2\left(\frac12-\frac2d\sum_jx_j\right)_+^2,\\
 g(x)&=S(x)+\frac2d\sum_j(x_j-5/8)^2,\qquad
 c(x)=\frac18+\frac1{8d}\sum_jx_j.
\end{align*}
The original scalar control has $0\le a\le c(x)$, running cost $S(x)+a^2+4a^4$, and transition
\[
 F_j(x,a,z)=\frac1{16}+\frac{x_j}{2}+\frac{x_{j+1}}8+
       \frac{x_j(1-x_{j+1})}{16}+B_ja+s_jz,
\]
where $(B_j,s_j)=(1/2,1)$ at alternating coordinates and $(1/4,-1)$ at the others, and $z$ is uniform on $[-1/32,1/32]$.

For the retained two-control extension, $a_1,a_2\ge0$, $a_1+a_2\le c(x)$, the second exposure column interchanges $1/2$ and $1/4$, and the action cost is
\[
 R(a)=a_1^2+a_2^2+4(a_1^4+a_2^4)+a_1a_2/4.
\]
An independent common-direction uniform innovation of radius $1/64$ is added. This is the previously specified extension, nested at $a_2=0$ when that innovation is disabled. Both models preserve $[0,1]^d$. Let $m\in\{1,2\}$ denote the number of controls and define
\begin{align*}
 Q_t^*(x,a)&=S(x)+R(a)+\beta\mathbb E V_{t+1}^*(F(x,a,Z)),\\
 V_t^*(x)&=\min_{a\in A(x)}Q_t^*(x,a),\qquad V_T^*=g.
\end{align*}
All costs and laws in this display are the economic primitives, not a fitted surrogate.

\begin{theorem}[Regularity of the original Bellman problem]\label{thm:regularity62}
For either control specification and every finite horizon, $V_t^*$ is convex and continuously differentiable in the interior, with the corresponding one-sided extension on the cube. It is $G_t$-Lipschitz in the $\ell^1$ norm and $M_t$-semiconcave in the Euclidean norm, where valid constants are
\begin{align*}
 G_T&=10/d,&G_t&=\frac{1037}{128d}+\beta\frac{47}{64}G_{t+1},\\
 M_T&=26/d,&M_t&=\frac{22}{d}+\frac{21}{256d}
             +\beta\left(\frac9{16}M_{t+1}+\frac18G_{t+1}\right).
\end{align*}
In particular, $G_t\le27/d$ and $M_t\le54/d$. The function $Q_t^*$ is jointly strongly convex on its convex feasible graph, with state curvature at least $107/(128d)$ and action curvature at least $7/4$.

A valid semiconcavity constant along each coordinate line is
\begin{align*}
 D_T&=\frac9d+\frac{16}{d^2},\\
 D_t&=\frac5d+\frac{16}{d^2}+\frac{21}{256d^2}
 +\beta M_{t+1}\left(\frac{85}{256}+\frac{23}{256d}\right).
\end{align*}
Consequently, for the positive multilinear interpolant $I_n$ on a mesh of spacing $1/n$,
\begin{equation}\label{eq:interpolation62}
 0\le I_nV_t^*(x)-V_t^*(x)\le\kappa_t,
 \qquad \kappa_t=\frac{dD_t}{8n^2}.
\end{equation}
\end{theorem}

\noindent The proof is in Supplement, Section~\ref{supp:proof-accuracy62-print-1}.


The theorem is a structural statement about these economic primitives. It is not a claim that every controlled diffusion, recursive preference specification, or game in the complete development edition has a convex Bellman problem. Nor does it require a signed-weight neural fit to preserve convexity. Convex approximation of dynamic programs has an established literature; the additional work here is to verify the needed curvature for the original nonlinear capacity-constrained model and connect it to the actual NBO actor and complete error account.

\subsection{Certified nodal Bellman bounds}
Suppose nodal arrays satisfy $l_{t+1}\le V_{t+1}^*\le u_{t+1}$ at all grid vertices. Equation~\eqref{eq:interpolation62} gives the functions $I_nl_{t+1}-\kappa_{t+1}$ and $I_nu_{t+1}$ as valid lower and upper continuations. The lower function need not itself be convex; convexity is used for the unknown true value and for the comparison argument, not asserted of arbitrary endpoint arrays.

Divide each uniform innovation component into $q$ equal bins and use its conditional mean. Convexity gives a lower quadrature bound for the true continuation. Semiconcavity gives an upper error
\begin{equation}\label{eq:quadrature62}
 \zeta_{t+1}=\frac{M_{t+1}}{6q^2}
                      \sum_{r=1}^{s}\rho_r^2\|s_r\|_2^2,
\end{equation}
where $s=1$ or $2$, $\rho_1=1/32$, and $\rho_2=1/64$ in the extension. The expectation remains the original continuous law. At the terminal date $g$ is evaluated directly, so no terminal interpolation allowance is needed.

At a vertex $v$, let $\mathcal A_A(v)$ be the capacity simplex lattice with spacing $c(v)/A$, including all its faces. For $m=1,2$, a minimizer can be represented by neighboring lattice vertices whose coordinate variances sum to at most $m c(v)^2/(4A^2)$. A valid action smoothness constant is
\[
 H_t=\frac{21}{4}+\beta M_{t+1}\|B\|_2^2,
 \quad \|B\|_2^2=\frac{5d}{32}\ (m=1),\quad
 \|B\|_2^2=\frac{9d}{32}\ (m=2).
\]
Taylor's inequality and that barycentric representation imply
\begin{equation}\label{eq:action-cover62}
 0\le\min_{a\in\mathcal A_A(v)}Q_t^*(v,a)-V_t^*(v)
       \le\delta_t:=\frac{mH_t}{128A^2}.
\end{equation}
The linear Taylor terms cancel even when the optimum lies on a face. This is an error against the continuous feasible optimum, not only against the finite action menu.

Let $q_t^-(v,a)$ and $q_t^+(v,a)$ include all interpolation, quadrature and arithmetic errors and enclose $Q_t^*(v,a)$. Then
\begin{equation}\label{eq:nodal62}
 l_t(v)=\max\{0,\min_{a\in\mathcal A_A(v)}q_t^-(v,a)-\delta_t\},
 \qquad
 u_t(v)=\min_{a\in\mathcal A_A(v)}q_t^+(v,a)
\end{equation}
are valid nodal Bellman bounds. The zero truncation uses the nonnegative economic costs. Extra neural or conventional proposals may improve the upper side but do not remove the lower action-cover allowance. Positive interpolation, averaging and minimum maps are monotone and nonexpansive; the separate allowances in \eqref{eq:interpolation62}--\eqref{eq:action-cover62} are therefore propagated, rather than being discarded because an endpoint calculation contains the truth.

\subsection{From a fitted witness to the complete implemented policy}
Let $a_t(v)\in A(v)$ be the actual action returned at each vertex by any generator. It may be selected by the exact neural witness, a conventional fitted continuation, or their common verified fallback. Write $w_v(x)$ for the multilinear weights and set $\bar\pi_t(x)=\sum_vw_v(x)a_t(v)$. Affineness of capacity makes this action feasible. Let $q_t^{\mathrm{sel},+}(v)$ be a verified upper bound on its original-optimum continuation value $Q_t^*(v,a_t(v))$, using the same lower and upper Bellman reference for all generators.

\begin{theorem}[Actual-policy, whole-domain certificate]\label{thm:policy62}
Set
\[
 e_t=\max_v\{q_t^{\mathrm{sel},+}(v)-l_t(v)\}+\kappa_t.
\]
Suppose the implemented policy $\widehat\pi_t$ is feasible and differs from $\bar\pi_t$ by at most $b_t$ in action $\ell^1$ norm, uniformly over the state cube. Define
\[
 \Lambda_t=\frac{13}{16}+\beta\frac{3d}{8}G_{t+1},\qquad
 D_T^\pi=0,\quad D_t^\pi=e_t+\Lambda_tb_t+\beta D_{t+1}^\pi.
\]
Then, at every true state and date,
\[
 0\le J_t^{\widehat\pi}(x)-V_t^*(x)\le D_t^\pi.
\]
This statement is conditional on the actual recorded policy, not on success of its fitting algorithm. It applies equally to a pure fitted policy and to a policy augmented by verified common candidates.
\end{theorem}

\noindent The proof is in Supplement, Section~\ref{supp:proof-accuracy62-print-2}.


In the execution, observation is rounded down at forty bits. An acquired state $y\le x$ has no greater capacity than the true state. Each interpolated action is then rounded down at twenty bits using the lower endpoint of its arithmetic enclosure. If $L_{ij}$ bounds the derivative of action coordinate $i$ with respect to state coordinate $j$, and $\eta$ bounds native interpolation error, a valid implementation allowance is
\[
 b_t\le 2^{-40}\sum_{i,j}L_{ij}
           +m(2^{-20}+2\eta+\eta_{\mathrm{out}}).
\]
Here $L_{ij}$ is computed from the maximum adjacent nodal difference times $n$, and $\eta_{\mathrm{out}}$ covers endpoint subtraction. All these numbers are recorded outward. This new barycentric acquired actor is explicitly distinguished from the earlier discontinuous cell-constant actor; an old actor is not silently replaced during a purported fixed-policy revalidation.

\subsection{Arithmetic, work, and scope}
The native interpolation kernel forms nonnegative tensor weights and a weighted sum of at most $2^d$ stored binary64 endpoints. It is compiled without reassociation or floating-point contraction. For $K=(3d+4)2^d+32$, an explicit conservative absolute allowance is
\[
 \eta=\frac{K2^{-53}}{1-K2^{-53}}\max_v|z_v|+K2^{-1022}.
\]
The final term conservatively covers underflow. Input coordinates, affine indexing, finite endpoints and the binary64 format are checked. Roundoff expansion and the ordinary interval cost arithmetic remain separate from the analytic approximation bounds. Exact rational fixtures compare the native sum with its mathematical tensor interpolant.

With $N=(n+1)^d$ vertices, $K_A=A+1$ scalar actions or $K_A=(A+1)(A+2)/2$ two-control actions, the reference calculation uses $O(TNK_Aq^s2^d d)$ elementary work under the explicit tensor implementation. Stored reference values use $O(TN)$ words and a policy uses $O(TNm)$ words. Streaming batches avoid a state-pair table. Proposal fitting, exact rational witness work, actor transfer, serialization, independent cost inference and compilation are additional costs, not absorbed into this bound. The analytic errors are second order in state, action and innovation resolution; the state-count and tensor-evaluation factors still depend exponentially on dimension. The result supplies a matched original-optimum comparison without asserting that these costs have disappeared.

\section{Neural Bellman queries without a state lattice}\label{sec:query63}
The preceding construction makes numerical error accountable, but a complete tensor reference can absorb the benefit of a fitted continuation. We therefore change the computational organization, not the economic problem. A Neural Bellman selector proposes an action at the state actually acquired. A recursive interval query then bounds the original Bellman value at that state and checks the proposed action against the entire continuous feasible interval. No conventional policy is precomputed at all states. The controller is the resulting callable algorithm, not the unverified proposal alone.

The distinction matters for attribution. Learning may reduce the work needed to locate a useful action. It does not supply the inequalities used to certify that action. The adaptive conventional comparator receives the same inequalities and the same original-optimum target. The trained value-fit policies in the preceding study remain distinct from the action-query predictor introduced here.

\subsection{A curvature lower bound on an adaptive action partition}
Use the original scalar economy of Section~\ref{sec:accuracy62}, with $r$ dates remaining. Write $V_r$ for its optimal value and $Q_r(x,a)$ for the value of a first action followed by the optimal future. Let $G_r,M_r$ be the primitive constants in Theorem~\ref{thm:regularity62}, indexed by remaining horizon, and put
\[
 H_r=\frac{21}{4}+\beta M_{r-1}\frac{5d}{32},
 \qquad \Lambda_r=\frac{13}{16}+\beta G_{r-1}\frac{3d}{8}.
\]
Thus $H_r$ is an upper curvature bound in the scalar action, and $\Lambda_r$ bounds the absolute action derivative. Neither depends on a fitted model.

\begin{lemma}[Parabolic minorants]\label{lem:parabolic63}
Suppose $f$ is $H$-semiconcave on $[l,u]$, and $L_l\le f(l)$, $L_u\le f(u)$. For $s\in[0,1]$,
\[
 f(l+s(u-l))\ge (1-s)L_l+sL_u
                 -\frac H2(u-l)^2s(1-s)=:p(s).
\]
Consequently $\min_{[0,1]}p$ is a lower bound for the minimum of $f$ on the interval. With $K=H(u-l)^2/2$ and $D=L_u-L_l$, this lower bound is $L_l$ if $D\ge K$, $L_u$ if $D\le-K$, and otherwise
\[
 \frac{L_l+L_u}{2}-\frac K4-\frac{D^2}{4K}.
\]
The endpoints $L_l,L_u$ need not themselves belong to a convex interpolant. Increasing $K$ preserves the lower-bound property.
\end{lemma}
\noindent The proof is in Supplement, Section~\ref{supp:proof-query63-1}.


At an acquired state, retain an ordered partition of the feasible action interval and an outward value bracket at every partition point. The minimum of the parabolic minorants bounds the continuous-action infimum from below. The smallest upper endpoint supplies an actual feasible witness. A neural or polynomial proposal inserts another point into this partition; it cannot remove the endpoints or any part of the feasible interval. The running lower bound is the maximum of all previously valid global lower bounds, so new information cannot weaken the certificate. The upper bound is the minimum over every evaluated feasible witness.

\subsection{Recursive queries under the original continuous law}
For an endpoint action with $r>1$, divide the scalar innovation into $q_r$ equiprobable bins and query $V_{r-1}$ at each conditional-mean next state. If each child returns width at most $\eta_r$, the endpoint bracket has width at most
\begin{equation}\label{eq:querywidth63}
 \beta\eta_r+\frac{\beta M_{r-1}d}{6144q_r^2}+\omega_r.
\end{equation}
Here $\omega_r$ includes outward primitive arithmetic and the Lipschitz expansion from an interval-valued next state to its queried center. Convexity provides the lower conditional-mean inequality; semiconcavity and the exact conditional variance provide the upper allowance. The innovation remains continuous. There is no state-interpolation term because a child value is queried rather than read from a state lattice.

For a requested local gap $\delta$, the implementation sets $\eta_r\le\delta/(4\beta)$ and chooses a dyadic $q_r$ so that the second term in \eqref{eq:querywidth63} is at most $\delta/4$. A separate endpoint-width guard charges arithmetic and the small capacity-rounding remainder. Failure of that guard is a precision failure, not a successful return. The terminal query uses analytic integration instead of recursive quadrature.

\begin{theorem}[A state-grid-free NBO service]\label{thm:query63}
Consider the original scalar investment economy. At every recursive query, suppose the outward endpoint and capacity-transfer accounts give width at most $5\delta/8$. Let the lower bound be the running maximum described above. Refine an interval attaining the smallest current minorant whenever the global upper-minus-lower gap exceeds $\delta$. Each new point must lie between one fifth and four fifths of that interval. Arbitrary additional learned proposals may be inserted initially. Then:

(i) every returned interval $[L_r(x),U_r(x)]$ contains $V_r(x)$, and its returned feasible action $a_r(x)$ satisfies
\[
 Q_r(x,a_r(x))-V_r(x)\le U_r(x)-L_r(x)\le\delta;
\]
(ii) the real-arithmetic refinement terminates. With at most one interior learned proposal, a sufficient point budget is
\begin{equation}\label{eq:querycap63}
 6+\left\lceil5c(x)\sqrt{H_r/(3\delta)}\right\rceil;
\end{equation}
(iii) for the policy that calls this service at every subsequent date, the original-law loss satisfies
\begin{equation}\label{eq:querypolicy63}
 J_T^{\pi}(x)-V_T(x)
 \le\sum_{t=0}^{T-1}\beta^t
 \left[\delta_t+
 \left(8+\frac{11\beta dG_{T-t-1}}{16}
                  +dG_{T-t}\right)2^{-b}\right],
\end{equation}
when the state is acquired downward at $b$ bits per coordinate and every returned action is feasible at that acquired state. All action rounding must precede its endpoint evaluation. The inequalities hold uniformly in the true initial state, not only at the states used in an experiment.
\end{theorem}
\noindent The proof is in Supplement, Section~\ref{supp:proof-query63-2}.


The finite-precision contract is explicit. The executed service uses the width guard, a 512-point cap, and a 32-bit downward action lattice. A precision or cap failure withholds the floating certificate and invokes a separately counted exact-rational enumeration. This fallback constructs values only at the required descendants; it does not load a precomputed conventional policy. The rational fallback chooses dyadic action and quadrature counts until each cover or quadrature allowance is at most one quarter of the request, and requests child width at most one quarter divided by the discount. Provided $\delta>8\Lambda_r2^{-32}$ at every recursive level, its exact brackets have width at most three quarters of the request. This proves a finite return for every acquired state in the declared tasks, including states absent from the workload. Rational bit-operation cost is additional work and is not replaced by a unit-cost floating clock. The mathematical point budget is checked against the declared tolerance range; termination in an arbitrary floating implementation is not inferred merely from the real-arithmetic argument. Fitting convergence is unnecessary for validity: any finite prediction, and even its clipped replacement when numerically invalid, is only a proposed query.

\subsection{A certified analytic terminal kernel}
The last decision admits an especially strong common comparator. Let $y=F(x,0,0)$, $b_0=1/2-2\operatorname{mean}(y)$, and $\gamma_j=B_j$. Define
\[
 D(x)=\frac8d\sum_j\gamma_j(y_j-5/8)
   +\frac1{2d}\sum_j(\gamma_j-\gamma_{j+1})(y_j-y_{j+1}).
\]
For every even $d$, the exact terminal action derivative is
\begin{equation}\label{eq:cubic63}
 16a^3+\left(2+\frac{41\beta}{32}\right)a
 +\beta D(x)-3\beta(b_0-3a/4)_+.
\end{equation}
Each branch is a strictly increasing cubic. Its floating closed-form root supplies only a candidate. The independently evaluated outward derivative and exact-law objective certify that candidate; a safeguarded derivative bracket remains available if the candidate is inaccurate.

For completeness, if $f$ is $\mu$-strongly convex and an outward interval encloses $f'(a)$, the largest compatible projected derivative residual $R$ gives
\[
 f(a)-\min_{[0,c]}f\le R^2/(2\mu).
\]
At zero, a nonnegative derivative has residual zero; at capacity, a nonpositive derivative has residual zero. Otherwise the absolute derivative bound is sufficient. This follows by minimizing the strong-convexity supporting quadratic over the feasible interval and, conservatively, over the line. Here $\mu=2$. Equation~\eqref{eq:cubic63}, its $41/32$ and $9/4$ coefficients, analytic shock moments, and outward derivatives are independently checked against rational formulas. Every method uses this kernel; it is not credited to the neural predictor.

\subsection{Work, prediction, and the economic object}
Let $A_r(\delta)$ be a bound on action endpoints at a query. Excluding training, a scalar operation account obeys
\[
 W_1(\delta)=O(dA_1),\qquad
 W_r(\delta)\le C d A_r(\delta)+
 A_r(\delta)q_r W_{r-1}(\delta/(4\beta)).
\]
The constants $dG_r,dM_r$ are bounded independently of $d$, so the declared endpoint and quadrature budgets need not grow exponentially with the state dimension. Primitive evaluations still cost $O(d)$. This removes the state lattice, not the horizon or shock-tree cost. The executed vectorization batches a query's innovation descendants. Its largest live state batch is reported; it does not claim the smaller memory of a fully streamed implementation. Increasing the horizon or the shock dimension may still be expensive. The new execution is scalar-control; the earlier full two-control result is retained rather than relabeled as a query-service experiment.

A useful prediction can reduce the endpoint count or return its own action from the initial partition. An inaccurate prediction can instead add work. This is a prediction-dependent algorithm, not a theorem of neural dominance. Learned warm starts have their own literature, including \citet{sakaue2023warm}; their discrete-convex setting differs from the original-law Bellman intervals here. Convex control policies and learned optimization parameters are studied by \citet{agrawal2020control}. Our comparison identifies the extra proposal cost and the resulting query transcript rather than claiming priority for learning a warm start.

The purchase is now a complete callable controller with a declared error contract. Its training cost, every subsequent certified query, and its return-time cost belong to the same service. A reduction in queried state storage is economically relevant for repeated decisions only after these costs are counted. The original cost normalization is maintained. Workload seconds and normalized expected-cost loss are reported separately; a price per unit of computation can convert the former to the latter only as an explicitly specified decision-maker parameter.

\section{Operational prediction certificates and readout training}\label{sec:operational67}
Proposition~\ref{prop:localization65} separates witness accuracy from lower-certificate localization. Its distance to an unknown minimizer is useful for analysis but is not a stopping observable. We now replace that distance by an upper envelope computed from information already purchased. This gives a certificate before an additional Bellman query, a training criterion measured in unresolved verification work, and a constructive neural readout with a computable optimization gap.

\subsection{A decision before the next query}
Fix a state, a date, and the current action partition. Write $f(a)=Q_t^*(x,a)$, and suppose $f$ is $\mu$-strongly convex on the feasible interval. Let $\overline f_l,\overline f_u$ be verified upper endpoints at two consecutive actions $l<u$. Let $L$ be the current lower bound for the original feasible infimum, including the capacity and arithmetic allowances. For $a=l+s(u-l)$, define
\begin{equation}\label{eq:screen67}
 C(a)=(1-s)\overline f_l+s\overline f_u
       -\frac\mu2 (a-l)(u-a),\qquad 0\le s\le1.
\end{equation}
Every symbol in this expression is either a primitive bound or an existing endpoint. No derivative or action at the unknown optimum occurs.

\begin{theorem}[Pre-query closure]\label{thm:screen67}
Under the preceding conditions, $f(a)\le C(a)$ throughout $[l,u]$. Let $a_\Delta\in[l,u]$ be the actual feasible rounded proposal. An outward upper endpoint $\overline C(a_\Delta)$ satisfying
\begin{equation}\label{eq:screen-stop67}
 \overline C(a_\Delta)-L\le\delta
\end{equation}
certifies $f(a_\Delta)-\inf_{a\in A(x)}f(a)\le\delta$ without another Bellman query. If the screen does not close, the protected refinement of Theorem~\ref{thm:gated64} retains its validity and finite query bound. The complete acquired policy therefore retains the original-optimum account of Theorem~\ref{thm:query63}.

Equivalently, the union over partition intervals of the sublevel sets $\{a:C(a)-L\le\delta\}$ is a computable certified action set. Each constituent is a quadratic sublevel interval intersected with its partition cell. Endpoints obtained numerically must be rounded inward, or the final action must be checked by \eqref{eq:screen-stop67}.
\end{theorem}
\noindent\emph{Proof.} See Supplement, Section~\ref{supp:proof70-operational67-1}.


The screen is shared by every comparator. In particular, minimizing the quadratic upper envelope is an available conventional proposal. Prediction is valuable only to the extent that its proposed location or evaluation cost improves the observed service. The executed neural rule screens its feasible prediction on the partition interval containing it, even when that location is outside the protected part of the currently critical interval. A failed screen does not authorize an unprotected split.

\subsection{A computable loss for unresolved work}
A \emph{certificate context} $z$ consists of the state, remaining horizon, tolerance, partition endpoints and their verified bounds. A proposal $a_\theta(z)$ is feasible and is computed without another Bellman query. Define its verified score and a margin loss by
\begin{align}
 s_\theta(z)&=\overline C_z(a_\theta(z))-L_z,\label{eq:score67}\\
 \ell_\tau(z,\theta)&=\min\left\{1,
       \frac{(s_\theta(z)-\delta_z+\tau)_+^2}{\tau^2}\right\},
       \qquad \tau>0.\label{eq:work-loss67}
\end{align}
Let $K(z)$ be a valid cap on further endpoint queries if the screen fails and the protected fallback is followed. This is a count of endpoint queries at the current state. The recursive cost of each endpoint is charged through the earlier descendant-work recursion. On contexts satisfying the declared precision contract, the real-arithmetic and rational-recovery bounds supply finite such caps.

\begin{theorem}[Certificate loss and additional work]\label{thm:work-loss67}
If $N_\theta(z)$ is the number of additional endpoint queries after the screen, then
\begin{equation}\label{eq:work-bound67}
 N_\theta(z)\le K(z)\mathbf1\{s_\theta(z)>\delta_z\}
                    \le K(z)\ell_\tau(z,\theta).
\end{equation}
Consequently, an achieved margin loss bounds the average additional queries on any specified finite context catalogue, with no unknown minimizing action.

For a distributional assertion, condition on construction of a finite family of $M$ predictors. Suppose $z_1,\ldots,z_n$ are independent validation contexts from a declared predictor-independent law, and $0\le K(z)\le K_{\max}$. With probability at least $1-\alpha$, simultaneously for that family,
\begin{equation}\label{eq:work-validation67}
 \mathbb E N_\theta(z)\le
 \frac1n\sum_{i=1}^n K(z_i)\ell_\tau(z_i,\theta)
 +K_{\max}\sqrt{\frac{\log(M/\alpha)}{2n}}.
\end{equation}
The cost of obtaining and validating contexts and of evaluating the predictor is additional work, not part of $N_\theta$.
\end{theorem}
\noindent\emph{Proof.} See Supplement, Section~\ref{supp:proof70-operational67-2}.


The context-law requirement matters. A stream of recursively generated contexts can depend on earlier predictions; it is not automatically an independent validation sample for \eqref{eq:work-validation67}. The deterministic inequality \eqref{eq:work-bound67} remains valid context by context. The recorded action-label mean squared error is not substituted for \eqref{eq:work-loss67}. The new loss measures a certificate's unresolved work rather than resemblance to an approximate action label.

\subsection{A constructive ReLU readout with an optimization certificate}
The work loss can be trained without invoking a successful hidden-layer optimizer. Fix a nonnegative affine--ReLU feature map $\phi(z)\in\mathbb R^p$ with $\sum_j\phi_j(z)=1$. Linear hat functions on a fixed coordinate partition are one exact construction; averages of such maps give multivariate features. A box-constrained readout $w\in[0,1]^p$ proposes
\[
 a_w(z)=l_z+(u_z-l_z)\phi(z)^\top w.
\]
The partition interval and features are fixed before optimizing $w$. Let $\eta_i\ge0$ cover action rounding and evaluation of the score, so that the computed score of the actual rounded action is at most $C_i(a_w(z_i))+\eta_i-L_i$. For example, on lattice-aligned endpoints an action Lipschitz bound gives $\eta_i=\Lambda_i\Delta_a$ for downward rounding. Define
\begin{equation}\label{eq:readout-loss67}
 F(w)=\frac1n\sum_i K_i
 \bigl(C_i(a_w(z_i))+\eta_i-L_i-\delta_i+\tau\bigr)_+^2
                      +\frac\lambda2\|w\|_2^2,
 \qquad \lambda\ge0.
\end{equation}
Unlike an action-label objective, this loss directly upper-bounds the additional-work account. It trains the upper witness against the current global lower certificate.

\begin{proposition}[Certified readout training]\label{prop:readout67}
The function $F$ is convex and continuously differentiable on the box. For $g=\nabla F(w)$, put
\begin{equation}\label{eq:dual-gap67}
 D(w)=g^\top w-\sum_{j=1}^p\min\{0,g_j\}.
\end{equation}
Then $0\le F(w)-\min_{v\in[0,1]^p}F(v)\le D(w)$. On the training contexts the average additional endpoint-query count is at most $F(w)/\tau^2$. If $\lambda>0$, the minimizer $w^*$ is unique and
\[
 \|w-w^*\|_2\le\sqrt{2D(w)/\lambda}.
\]
All quantities in these conclusions can be bounded by outward arithmetic; they are exact rational quantities when the contexts, features and iterate are rational.
\end{proposition}
\noindent\emph{Proof.} See Supplement, Section~\ref{supp:proof70-operational67-3}.


This is a fixed-feature neural construction, not a convergence assertion for unrestricted nonconvex network fitting. The deposited rational implementation uses finite conditional-gradient steps and reports its achieved loss and dual gap. Its regression examples test the proposition; they are not new investment-policy observations. The preceding R67 service comparison retained its preregistered action-label fits. The new comparison in Section~\ref{sec:study70} separately executes this certificate-aware objective, with equal-information Bernstein and label-loss controls; no historical fit is relabeled as its minimizer. Prediction-dependent optimization has an established literature \citep{sakaue2023warm,sambharya2023}; the additional object here is a computable original-law Bellman certificate that can license an action before another query and supplies a directly auditable training criterion.

\section{Two controls and two continuous innovations}\label{sec:vector67}
The absence of a state lattice does not remove action or innovation dimensions. We extend the query implementation to the previously specified two-control economy of Section~\ref{sec:accuracy62}: $a_1,a_2\ge0$, $a_1+a_2\le c(x)$, the interaction cost is $a_1a_2/4$, and the second innovation is an independent common-direction uniform shock. The original scalar model and its experiments are unchanged.

At a queried state, triangulate the feasible action simplex. Let $v_0,v_1,v_2$ be a triangle's vertices with verified lower endpoints $l_0,l_1,l_2$. If $\theta_j(a)$ are its barycentric coordinates, define
\begin{equation}\label{eq:triangle67}
 p(a)=\sum_{j=0}^2\theta_j(a)l_j
       -\frac H2\sum_{j=0}^2\theta_j(a)\|v_j-a\|_2^2.
\end{equation}
This is the quadratic $H\|a\|_2^2/2+b^\top a+c$ interpolating the lower endpoints. Its approximate minimizer is only a proposed support point. For any point $y\in\mathbb R^2$, including a numerically displaced point outside the triangle,
\begin{equation}\label{eq:dual-triangle67}
 p(y)+\min_j\nabla p(y)^\top(v_j-y)
                    \le \min_{a\in\operatorname{co}(v_0,v_1,v_2)}p(a).
\end{equation}
Thus an outward supporting-plane calculation, not an uncertified quadratic solver, supplies the lower endpoint.

\begin{theorem}[Simplicial Bellman query]\label{thm:vector67}
For the two-control primitives, \eqref{eq:triangle67} is a lower minorant of the true Bellman objective. Taking the smallest valid triangle lower endpoint and the best feasible-action upper endpoint gives a lower--upper certificate for the full continuous simplex. Include the omitted capacity strip and actual action-rounding allowances. Returning when this gap is at most $\delta$ gives the same acquired-policy loss recursion as Theorem~\ref{thm:query63}.

A failed floating triangulation can be replaced by an exact-rational simplex enumeration. Under the stated positive precision slack, this recovery is finite at every finite horizon. With two innovations divided into $q$ bins each, a valid local quadrature error is
\begin{equation}\label{eq:two-shock67}
 \zeta=\frac{5\beta M_{t+1}d}{24576q^2}.
\end{equation}
If $A_r$ bounds the number of action queries at a state with $r$ dates remaining, a work account has the form
\begin{equation}\label{eq:vector-work67}
 W_r\le C_d A_r+A_rq_r^2W_{r-1},\qquad C_d=O(d),
\end{equation}
with separate charges for triangulation, arithmetic and output. There is no tensor state-table factor, but the descendant tree remains.
\end{theorem}
\noindent\emph{Proof.} See Supplement, Section~\ref{supp:proof70-vector67-1}.


The executed algorithm uses longest-edge refinement of the lowest triangle and an outward terminal convexity-gap check. Its floating cap is followed by the stated recovery rule, not by an uncertified returned action. The implementation batches descendants but can allocate a descendant array before slicing it into query batches; a reported maximum query batch is therefore not a bound on total live state memory. We report process memory separately. The two-control execution is a verified extension of the NBO query backend, not evidence for a trained two-output selector.

\section{Frozen neural continuations and policy-cost precision}\label{sec:frozen72}
The optimal-advantage estimator in Section~\ref{sec:centered69} makes a policy comparison informative, but its conditional optimal-value queries are expensive. A fitted continuation can replace those queries without changing the estimand. The necessary certificate concerns its Bellman residual across states, not its distance from an unknown optimizer. This places the original NBO continuation, rather than only a proposed action, inside the evaluation calculation.

\subsection{Residual oscillation is the relevant range}
Consider the bounded expected-cost model with common terminal cost $g$, discounts $\beta_t\ge0$, and a common initial-state law. Write $b_0=1$ and $b_t=\prod_{j<t}\beta_j$. Fix bounded measurable continuations $h_0,\ldots,h_T$, conditionally on their construction. Define
\begin{align*}
 Q_t^h(x,a)&=c_t(x,a)+\beta_tP_t^a h_{t+1}(x),\\
 r_t^h(x)&=\inf_{a\in A_t(x)}Q_t^h(x,a)-h_t(x),\\
 a_t^{h,\pi}(x)&=Q_t^h(x,\pi_t(x))-\inf_{a\in A_t(x)}Q_t^h(x,a).
\end{align*}
The infimum is over the original continuous feasible set. Suppose verified constants give $0\le a_t^{h,\pi}\le\eta_t^\pi$, and likewise for $\nu$, throughout the true state domain. Let $\omega_t$ bound the oscillation of $r_t^h$, and $\omega_T$ that of $g-h_T$. Set
\[
 \Omega_h=\sum_{t<T}b_t\omega_t+b_T\omega_T,
 \qquad H_h^\pi=\sum_{t<T}b_t\eta_t^\pi.
\]
These constants include acquisition, action rounding and arithmetic allowances wherever those operations affect the implemented policy or the verified functions. A small empirical residual alone is not such a certificate.

\begin{theorem}[Frozen-continuation comparison]\label{thm:frozen72}
Along the actual trajectory of $\pi$, put
\begin{equation}\label{eq:frozen-observation72}
 Z_h^\pi=\sum_{t<T}b_t\bigl[Q_t^h(X_t^\pi,\pi_t(X_t^\pi))-h_t(X_t^\pi)\bigr]
                 +b_T(g-h_T)(X_T^\pi).
\end{equation}
For any coupling of the two original trajectory laws, $D_h=Z_h^\pi-Z_h^\nu$ satisfies
\begin{align}
 \mathbb E D_h&=\mathbb E[J_0^\pi-J_0^\nu],\label{eq:frozen-identity72}\\
 -\Omega_h-H_h^\nu&\le D_h\le\Omega_h+H_h^\pi.\label{eq:frozen-range72}
\end{align}
Moreover, at every initial state,
\begin{equation}\label{eq:frozen-regret72}
 0\le J_0^\pi-V_0^*\le\Omega_h+H_h^\pi.
\end{equation}
Adding an arbitrary state-independent constant to each $h_t$ changes neither $D_h$ nor these oscillation bounds. With $h_t=V_t^*$, the result reduces to optimal-advantage centering.
\end{theorem}
\noindent The proof is in Supplement, Section~\ref{supp:frozen72}.

The terminal correction in \eqref{eq:frozen-observation72} is essential unless $h_T=g$. Conditional expectation, not pathwise equality with the realized cash-cost difference, supplies \eqref{eq:frozen-identity72}. The result uses the actual kernels and costs; it assumes neither a correct statistical transition model estimated from data nor a convex neural fit. The pointwise policy certificate also applies conditionally after an adaptive choice among valid continuations. The independent-sample confidence statement below requires its additional sampling conditions.

For $M$ conditionally fixed comparisons, suppose independent trajectory rows yield outward enclosures $[L_i,U_i]$ of $D_h$, intersected with its valid support. Write $R_h=2\Omega_h+H_h^\pi+H_h^\nu$. A simultaneous interval with failure probability at most $\alpha$ is
\begin{equation}\label{eq:frozen-ci72}
 \left[\overline L-R_h\sqrt{\frac{\log(2M/\alpha)}{2n}},\quad
       \overline U+R_h\sqrt{\frac{\log(2M/\alpha)}{2n}}\right].
\end{equation}
This follows by the bounded-sum inequality applied to both endpoint samples and a union bound. The numerical and sampling width allocations of Corollary~\ref{cor:precision69} apply with $R_h$. Fitting, uniform residual qualification, conditional integration and any subsequent selection are distinct paid operations. In particular, a cache that adapts across rows changes the controller's history: its rows cannot be treated as independent evaluations of one frozen Markov policy. A fixed cache or an independent episode-reset design is sufficient to retain the stated sampling construction.

\subsection{Conditional integration of an actual neural continuation}
A shallow rational affine--ReLU continuation has the form
\[
 h(y)=c+\ell^\top y+\sum_{j=1}^K v_j(w_j^\top y+u_j)_+.
\]
Its readout weights may have either sign. In the original investment laws the next state, conditional on $(x,a)$, is affine in one or two independent centered uniform innovations. A neuron therefore has projected center $z$ and nonnegative projection radii $A,B$. Write $z_+=\max\{0,z\}$.

\begin{proposition}[Exact conditional ReLU moments]\label{prop:moment72}
For independent uniform projections of radii $A,B>0$, the conditional means of a neuron are
\begin{align}
 \mathcal M_1(z;A)&=\frac{(z+A)_+^2-(z-A)_+^2}{4A},\label{eq:moment-one72}\\
 \mathcal M_2(z;A,B)&=\frac{1}{24AB}\bigl[(z+A+B)_+^3-(z+A-B)_+^3\nonumber\\
 &\hspace{37mm}-(z-A+B)_+^3+(z-A-B)_+^3\bigr].\label{eq:moment-two72}
\end{align}
A zero projection radius is removed; if all radii vanish the mean is $z_+$. Consequently $P^ah$ is a finite exact rational calculation for rational primitives, states, actions and weights. For $n$ trajectories and at most $K$ neurons per date, its arithmetic-operation count is $O(nTKd)$, exclusive of fitting, trajectory construction and uniform qualification. No optimal-value descendant query occurs in this calculation.
\end{proposition}
\noindent The proof is in Supplement, Section~\ref{supp:moments72}.

The operation count does not bound rational bit complexity. Near-zero projected radii can also make a naive floating evaluation unstable; directed rounding or rational evaluation, including the zero-radius branches, is required. The executed moment catalogue checks the rational identities against an independently written polygon integral. It is a conditional-integration experiment, not a new trained-policy cost comparison. The formulas do not imply that the achieved residual oscillation of an arbitrary fitted network is small.

\subsection{Qualifying a fitted continuation without a Bellman tree}
Uniform qualification need not query $V^*$. Suppose $Q_t^h$ is Lipschitz in state and action with constants $L_{x,t},L_{a,t}$, and $h_t$ has constant $L_{h,t}$, all in $\ell^1$. In the investment capacity simplex, an action feasible at $y$ can be repaired continuously at $x$ with displacement at most $\kappa\|x-y\|_1$, where $\kappa=1/(8d)$. Thus a valid residual Lipschitz constant is
\[
 L_{r,t}=L_{x,t}+\kappa L_{a,t}+L_{h,t}.
\]
Let a finite state cover have fill radius $\rho$, and let $[l_j,u_j]$ enclose $r_t^h$ at every cover point. Then
\begin{equation}\label{eq:residual-cover72}
 \min_jl_j-L_{r,t}\rho\le r_t^h(x)
       \le\max_ju_j+L_{r,t}\rho
\end{equation}
throughout the cube. At a cover point, a feasible action mesh of fill radius $s$ gives a lower infimum bound by subtracting $L_{a,t}s$ from the smallest action-value lower endpoint; the smallest upper endpoint is an upper bound. Exact moments can supply those action endpoints. None of these steps requires convexity of $h$ or a true-value curvature minorant for $Q^h$.

For $N_t$ state-cover points and $A_t$ actions per point, this qualification costs $O(\sum_tN_tA_tK_{t+1}d)$ arithmetic operations, plus certification of the moduli and any actor-specific greedy-gap bound. Its cover can still be expensive in high dimension. The result replaces a recursive optimal-value oracle by an explicit approximation-and-qualification tradeoff; it does not make that qualification free. The residual-cover proof, the action-mesh argument and the terminal allowance are given in the supplement. Value-function control variates have an established variance-reduction literature \citep{munos2006variance}; the additional requirement here is an original-law, whole-domain certificate for the fitted NBO continuation and the actual acquired policy.

\section{Fitting a continuation for qualification}\label{sec:qualification73}
A continuation used to evaluate a policy should be fitted to the error that determines the evaluation cost. The residual-cover bound in \eqref{eq:residual-cover72} supplies such an objective. We now make its optimization constructive. The future continuation is held fixed at each backward step; neither an optimal-value oracle nor success of unrestricted hidden-layer training is assumed.

Fix a date and suppress its subscript. Let $x_1,\ldots,x_N$ be a state cover of radius $\rho$ in the $\ell^1$ metric, and suppose
\[
 y_j^-\le (\mathcal T h_+)(x_j)\le y_j^+.
\]
These are verified bounds on the original continuous-action infimum, not only on a proposal menu. Write $h_\theta=\sum_{k=1}^K\theta_k\phi_k$, where the fixed features have verified Lipschitz constants $L_k$ and $|\theta_k|\le B_k$. The constant feature is allowed with $L_k=0$. Let $C$ be a verified Lipschitz constant of $\mathcal T h_+$, including state-dependent feasibility repair. With $F_{jk}=\phi_k(x_j)$, define
\begin{equation}\label{eq:readout-objective73}
 \mathcal W(\theta)=\max_j(y_j^+-F_j\theta)
       -\min_j(y_j^--F_j\theta)
       +2\rho\left(C+\sum_k L_k|\theta_k|\right).
\end{equation}

\begin{theorem}[Qualification-optimal readout]\label{thm:readout73}
For every feasible $\theta$, $\operatorname{osc}(\mathcal T h_+-h_\theta)\le\mathcal W(\theta)$. Minimizing \eqref{eq:readout-objective73} is a finite linear program with $2K+2$ variables and $2N+2K$ inequality constraints, apart from variable bounds. For probability vectors $p,q$ on the nodes, set $\lambda_k=2\rho L_k$ and
\begin{equation}\label{eq:readout-dual73}
 \mathcal D(p,q)=p^\top y^+-q^\top y^-+2\rho C
   -\sum_k B_k\bigl(|[F^\top(q-p)]_k|-\lambda_k\bigr)_+.
\end{equation}
Then $\mathcal D(p,q)\le\inf_\theta\mathcal W(\theta)$, with equality for some $p,q$. Thus any feasible readout and exact probability vectors certify an achieved optimization gap of at most $\mathcal W(\theta)-\mathcal D(p,q)$.
\end{theorem}

The proof is in the Technical Supplement. It constructs the primal constraints and dual support function explicitly. Rationalizing a solver's proposed readout and probability vectors gives a checkable certificate; a convergence flag is not the certificate. Label intervals, features, moduli, coefficient bounds and the cover remain part of the numerical contract.

Applying the theorem backwards gives a fitted continuation at every date. Its achieved ranges $\mathcal W_t$, together with the terminal oscillation and the actual implemented greedy-gap bounds, can be substituted directly into Theorem~\ref{thm:frozen72}. This connects the attained fitting objective to both original-optimum regret and the range of the original-law cost observation. Each readout minimizes its stated one-step qualification bound conditional on the fixed future. It does not minimize the joint multi-date range, the true policy loss, or a hidden-feature objective. The next backward step recomputes its labels and moduli from the continuation actually returned.

The label construction in \eqref{eq:residual-cover72} uses exact conditional moments for a shallow ReLU future and a verified continuous-action cover. Its work is charged in addition to the linear program, coefficient rationalization, modulus calculation and dual check. A conventional basis with the same number of coefficients receives the same labels and the same objective; it uses its own valid feature moduli. This is a comparison of qualified approximations, not an assertion that different bases span the same functions or that a neural encoding has an intrinsic advantage.

\section{Sequential qualification of policy costs}\label{sec:sequential73}
A fixed-sample interval does not by itself justify repeated inspection until a controller appears acceptable. Nor do cost observations from changing controllers identify the last controller's cost. We extend the frozen-continuation identity to predictable episodes and distinguish these targets. This uses established time-uniform concentration principles \citep{howard2020time}; confidence sequences have also been used for gated deployment in contextual bandits \citep{karampatziakis2021}. Here the observations are verified, continuous-law NBO cost contrasts, not importance-weighted bandit rewards.

Let $\mathcal F_{i-1}$ contain all information before episode $i$. Conditional on it, choose two complete controllers, a common initial law, a continuation $h_i$ fixed through the episode, and valid coefficients of Theorem~\ref{thm:frozen72}. The resulting cost contrast $D_i$ has conditional mean $\Delta_i$ equal to the original-law expected cost difference and predictable support $[a_i,b_i]$. Let $[L_i,U_i]\subseteq[a_i,b_i]$ enclose $D_i$ for the entire continuous innovation cell. Write $V_n=\sum_{i\le n}(b_i-a_i)^2$. For fixed positive numbers $\lambda_j,w_j$, $1\le j\le J$, with $\sum_jw_j\le1$, define
\begin{equation}\label{eq:anytime-boundary73}
 B_n=\min_{1\le j\le J}
       \left\{\frac{\log(2/(\alpha w_j))}{\lambda_j}
                    +\frac{\lambda_j V_n}{8}\right\}.
\end{equation}

\begin{theorem}[Predictable original-law cost certificate]\label{thm:anytime73}
Under the preceding conditional-law, support and enclosure conditions, with probability at least $1-\alpha$, simultaneously for every $n\ge1$,
\begin{equation}\label{eq:anytime-interval73}
 \frac{\sum_{i\le n}L_i-B_n}{n}
 \ \le\ \frac1n\sum_{i\le n}\Delta_i
 \ \le\ \frac{\sum_{i\le n}U_i+B_n}{n}.
\end{equation}
Each endpoint may also be intersected with the corresponding average predictable support endpoint. The assertion holds at every almost surely finite stopping time. If the complete controller pair and initial-law/reset contract stay fixed, then $\Delta_i=\Delta$, even when the qualified continuation changes predictably between episodes. Otherwise the target is the average conditional cost difference, not the last controller's cost.
\end{theorem}

The proof applies conditional Hoeffding bounds to each endpoint and Ville's inequality to a finite collection of nonnegative supermartingales. It requires no independent or identically distributed deployment contexts. It does require the stated conditional transition laws. Within-episode cache adaptation is allowed only when the controllers are treated as history-dependent algorithms and the greedy-gap qualification covers every memory state that can occur. Repeating the same program identifier while carrying a different cache state does not establish a fixed conditional target.

The implementation uses a finite dyadic collection of $\lambda_j$ and exact rational upper bounds for the logarithms. These give conservative, computable versions of \eqref{eq:anytime-boundary73}. It reserves the support and target identity before accepting an episode. A failed numerical evaluation consumes the full previously qualified support interval and its actual work; it is not deleted after inspecting its outcome. A failure of the support certificate itself stops qualification.

\begin{corollary}[Costed replacement after optional stopping]\label{cor:adoption73}
For a fixed target pair, let $[l_N,u_N]$ be the interval at a finite stopping time $N$. Let a declared deployment count be $m\ge1$, a one-time replacement and qualification charge be $F\ge0$, and the incremental per-deployment charge be $c$, all in the same units as the policy cost. On the simultaneous event of Theorem~\ref{thm:anytime73}, the net incremental cost lies in
\[
 [l_N+F/m+c,\ u_N+F/m+c].
\]
Consequently $u_N+F/m+c<0$ certifies strict net improvement. Inclusion of $[l_N,u_N]$ in a declared equivalence band certifies equivalence at that band, not strict superiority. New adaptively selected controller pairs may be qualified in separate epochs with conditional error allocations whose sum is at most the overall error budget.
\end{corollary}

A wall-clock reading is not a monetary charge. Measured setup, validation, rejected evaluations and successful retrievals are retained separately; prices are explicit decision inputs. The number of qualification episodes is not automatically the number of deployments over which a setup cost is amortized. These distinctions let the original investment decision use a valid stopping rule without silently changing either its estimand or its units.

\section{Original-law qualification of constructed policies}\label{sec:cost77}
A continuation can have a verified fitting objective without its implemented policy having an identified economic advantage. The distinction is resolved by evaluating the same complete policy that the construction returns. We now connect the qualification-optimal readouts of Section~\ref{sec:qualification73} to independent original-law cost observations. Neither the fitted future nor the acquired actor is replaced during evaluation. The comparison concerns the constructed NBO policies, rather than a different recursive controller supplied with the same neural label.

\subsection{Conditional integration and numerical observations}
In the two-state, two-date investment economy, write $a_t=\pi_t(X_t)$ and $\beta=15/16$. A cash observation and a conditionally integrated observation are
\begin{align}
 C^\pi&=c(X_0,a_0)+\beta c(X_1,a_1)+\beta^2g(X_2),\label{eq:cash77}\\
 Y^\pi&=c(X_0,a_0)+\beta c(X_1,a_1)
                    +\beta^2\bar g(X_1,a_1),\label{eq:conditional77}
\end{align}
where $\bar g(x,a)=P^ag(x)$. If $m(x,a)$ is the deterministic part of the transition, the original common, opposite-signed innovation gives the exact identity
\begin{equation}\label{eq:last-moment77}
 \bar g(x,a)=g(m(x,a))+\frac{5}{3072}.
\end{equation}
The shortage term depends on the sum of the coordinates and is unchanged by that innovation. The quadratic state and difference terms supply the variance correction. Consequently,
\[
 \E Y^\pi=\E C^\pi=\E J_0^\pi(X_0),\qquad
 \operatorname{Var}(Y^\pi-Y^\sigma)
       \le\operatorname{Var}(C^\pi-C^\sigma)
\]
when both policies use the same initial state and innovations. The variance inequality concerns the exact observations; it does not assert that every finite-sample numerical interval must be narrower.

The state cube is invariant. On that cube, $0\le a\le1/4$, a stage cost is at most $103/64$, and the terminal cost is at most $37/16$. Both observations therefore lie in
\begin{equation}\label{eq:support77}
 [0,H],\qquad H=\frac{103}{64}(1+\beta)+\frac{37}{16}\beta^2
                  =\frac{21097}{4096}.
\end{equation}
A paired difference has the signed support $[-H,H]$. Negative common-path gains are not clipped. No original-optimum value query is required to compute either observation for a frozen acquired actor.

An initial or innovation bin is a continuous cell, not a point mass. The evaluator encloses every true state in the bin, every acquired cell reached by that state enclosure, and every resulting feasible action. When an enclosure crosses an acquisition boundary, all intersected actions enter the continuation enclosure. A second implementation evaluates the same primitives using integer-directed arithmetic and expanded polynomials. The reported endpoints are the hull of the two valid enclosures, intersected with the proved support. Arithmetic agreement is a check on the implementations; validity follows from the enclosure rules and the primitive identities, not from agreement alone.

\subsection{Precision and the decision criterion}
The following statement applies to any fixed finite collection of qualified policies, not only to the two feature bases used below. All construction information is fixed before the independent evaluation rows are drawn. Dependence among policies on a common row is allowed.

\begin{theorem}[Original-law cost qualification]\label{thm:precision77}
For each of $K$ declared estimands, let $Z_i$ be its original-law observation on independent row $i$, and suppose a deterministic full-bin calculation returns
\[
 a\le L_i\le Z_i\le U_i\le b,\qquad i=1,\ldots,N.
\]
The inclusion may instead hold for the original conditional expectation within the sampled bin. Put $R=b-a$, let $\bar L,\bar U$ denote sample means, and let $s_L^2,s_U^2$ be the unbiased endpoint sample variances. For $N\ge2$ and $\ell>0$, define
\[
 r_L=\sqrt{\frac{2s_L^2\ell}{N}}+\frac{7R\ell}{3(N-1)},\qquad
 r_U=\sqrt{\frac{2s_U^2\ell}{N}}+\frac{7R\ell}{3(N-1)}.
\]
With probability at least $1-4K e^{-\ell}$, simultaneously for the declared collection,
\begin{equation}\label{eq:interval77}
 \E Z_i\in[\bar L-r_L,\bar U+r_U].
\end{equation}
Before rounding the final endpoints, its width is exactly
\begin{equation}\label{eq:width77}
 (\bar U-\bar L)
 +\sqrt{\frac{2s_L^2\ell}{N}}+\sqrt{\frac{2s_U^2\ell}{N}}
 +\frac{14R\ell}{3(N-1)}.
\end{equation}
Upward square roots and outward final rounding preserve coverage. The first term includes all retained numerical and acquisition uncertainty; the middle terms reflect observed endpoint dispersion; the last term is the finite-support penalty. A declared full-width target is met only when their computed sum meets it.
\end{theorem}

The proof in the Technical Supplement applies the empirical-Bernstein inequality of \citet{maurer2009} to the lower and upper endpoint variables. The statistical inequality is established work. The contribution here is its application to the original-law costs of the actual constructed actors, with executable whole-bin observations, independent arithmetic and an explicit precision account. Refining only a simulation bin cannot remove the finite-support or sampling terms. Conversely, an inaccurate uniform regret certificate need not prevent an accurate comparison of two actual policy costs.

Precision, equivalence and improvement are different conclusions. For a cost-difference interval $[l,u]$ and a predeclared margin $\delta$, equivalence requires $-\delta<l\le u<\delta$. Noninferiority of the first policy requires $u<\delta$; strict superiority requires $u<0$. An interval can meet its width target and fail equivalence because it identifies a material difference. These decisions are made from exact endpoints, not from rounded table entries.

\begin{corollary}[Priced selection and adoption]\label{cor:purchase77}
On the simultaneous event of Theorem~\ref{thm:precision77}, let $[L_j,U_j]$ qualify the cost of policy $j$, and let $w_j$ be its declared nonnegative resource charge. A minimizer $\widehat j$ of $U_j+w_j$ satisfies, for every policy $j$,
\[
 J_{\widehat j}+w_{\widehat j}
       \le J_j+w_j+(U_j-L_j).
\]
The inequality holds for every assignment of resource prices on the same event. If the gain from replacing an installed controller is in $[l,u]$, then $m$ deployments, monetary scale $v\ge0$, qualification charge $q\ge0$ and replacement fee $f\ge0$ give expected net benefit in
\[
 [mv l-q-f,\;mv u-q-f].
\]
In particular, when $l>0$ and $v>0$, every $m>(q+f)/(vl)$ is certified beneficial.
\end{corollary}
\begin{proof}
The first assertion follows from $J_{\widehat j}\le U_{\widehat j}$, the minimizing property, and $U_j\le J_j+U_j-L_j$. These deterministic inequalities hold for all nonnegative prices once all cost intervals cover. The second assertion is multiplication by $mv$ and subtraction of $q+f$.
\end{proof}

The resource charges in this result must be specified, not inferred from normalized losses. An already paid construction and a prospective purchase have different charges. When only a shared qualification service has been timed, charging that entire service to a proposed replacement is conservative; inventing a per-policy marginal qualification time is not.

\subsection{The constructed-controller experiment}
The fixed catalogue consists of the four previously constructed NBO actors and the admissible zero-investment incumbent. ReLU and degree-four Bernstein continuations each have nine coefficients and are constructed on eight- or sixteen-cell coordinate covers. At a fixed cover, the two bases use the same acquisition contract and primitive information. The actor arrays, actual fitted futures and exact optimization certificates are unchanged. The zero-investment policy supplies an economically interpretable installed reference; it is not substituted for the strong conventional query competitors in the preceding study.

Before execution, the protocol fixes $N=524{,}288$, forty-bit continuous bins, the sample stream, the full-width target $1/256$, and the separate two-sided equivalence margin $1/256$. All five marginal costs and ten paired differences are reported for both observations in \eqref{eq:cash77}--\eqref{eq:conditional77}. Thus $K=30$ and $\ell=10$ give a joint failure bound below $1/100$. Reusing a row for another controller or estimator does not increase the independent sample count. The declared independence model concerns the simulated rows; a fixed pseudorandom seed establishes reproducibility, not mathematical independence.

\begin{table}[htbp]\centering
\caption{Original-law costs of the constructed controllers}\label{tab:cost77}
{\small\setlength{\tabcolsep}{4pt}\begin{tabular}{lcc}
\toprule
Controller & Expected cost & Gain over zero incumbent \\
\midrule
Zero & [0.773561, 0.780120] & -- \\
R8 & [0.607716, 0.613107] & [0.165305, 0.167552] \\
B8 & [0.684034, 0.690159] & [0.089044, 0.090443] \\
R16 & [0.607022, 0.612384] & [0.165995, 0.168280] \\
B16 & [0.682910, 0.689022] & [0.090168, 0.091580] \\
\bottomrule\end{tabular}}
\par\smallskip{\footnotesize R and B denote the original nine-coefficient ReLU and Bernstein constructions; the number is the coordinate cover size. All policies are frozen before the new evaluation rows. The displayed observations integrate the last innovation analytically without changing expected cost. Endpoints are rounded outward. Paired gain intervals, not differences of marginal cost intervals, determine gain signs.}
\end{table}

\begin{table}[htbp]\centering
\caption{Paired feature and acquisition comparisons}\label{tab:pairs77}
{\small\setlength{\tabcolsep}{4pt}\begin{tabular}{lcrc}
\toprule
First minus second & Expected-cost difference & Full width & Equivalent \\
\midrule
R8 $-$ B8 & [-0.077568, -0.075801] & 0.001767 & No \\
R16 $-$ B16 & [-0.077159, -0.075368] & 0.001791 & No \\
R8 $-$ R16 & [0.000227, 0.001191] & 0.000963 & Yes \\
B8 $-$ B16 & [0.000661, 0.001599] & 0.000937 & Yes \\
\bottomrule\end{tabular}}
\par\smallskip{\footnotesize A negative interval identifies lower cost for the first controller. Equivalence is determined from exact endpoints at the predeclared two-sided margin 1/256. The separate full-width target is also 1/256. Matched-basis comparisons hold the acquisition contract fixed; cross-resolution comparisons do not.}
\end{table}

\paragraph{Outcomes of the fixed precision design.}

Across both estimators, 20 of the twenty paired intervals meet the declared full-width target. Their widths range from 0.000936 to 0.002285. Exactly 4 meet the separate equivalence criterion. Neither a failed width target nor a failed equivalence comparison is omitted.

The conditionally integrated comparisons identify a positive gain over the zero-investment incumbent for 4 of the four constructed policies. These are gains of complete implemented controllers, not reductions in a qualification objective.

At the two matched acquisition resolutions, the ReLU-minus-Bernstein comparisons identify 2 lower-cost ReLU outcomes, 0 higher-cost outcomes, and 0 unresolved signs. This finite comparison is not a universal representation ranking.


The construction and the evaluation now refer to the same complete controllers. The resulting cost intervals, rather than a reduction in the residual-fitting objective, determine the adoption regions in Corollary~\ref{cor:purchase77}. Comparisons between the feature bases identify only what their paired intervals support. The earlier unresolved and adverse query-method comparisons remain distinct evidence and are preserved in the Numerical and Development Appendix.

Both interval implementations evaluate every new trajectory and every acquired-action rectangle. Additional rational midpoint checks test predetermined rows. The phase ledger separates controller loading, lookup preparation, binary64 evaluation, integer evaluation, exact statistical arithmetic, output and the complete parent process. The independent replay recomputes the entire new enclosure catalogue without replacing its original clocks. This evaluates the two-date constructed-policy problem; it is not an independent rerun of the older, much larger vector Bellman tree.

\section{Deployment choice under qualified cost uncertainty}\label{sec:allocation79}
The attained continuation is an input to a complete economic decision, not merely to another approximation score. The same qualified costs can support a choice of acquisition precision across deployments. We retain the original NBO actors and draw a controller before observing the initial state or innovations. This implements a mixture of complete policies, not an average of their actions, coefficients or information sets. Its expected original-law cost is the corresponding mixture of the qualified costs. Minimax-regret allocation has an established decision-theoretic interpretation \citep{manski2004}; the result below concerns the finite cost set supplied by NBO qualification, conditional on its simultaneous coverage event.

Let $\mathcal C$ and $d(j,i)$ be the consistent difference-constraint set and distances of Theorem~\ref{thm:choice78}, reproduced with its full proof in Appendix~\ref{app:choice78}. For a probability vector $p\in\Delta_q$ and per-deployment charges $w$, set $H_{ji}(w)=d(j,i)+w_i-w_j$, for $i,j=1,\ldots,q$. The charges include all costs that actually vary linearly with deployment shares; additional installation costs are treated separately.

\begin{theorem}[Sharp deployment allocation]\label{thm:allocation79}
The largest expected regret of allocating complete controllers with probabilities $p$ is
\begin{equation}\label{eq:allocation79}
 \mathcal R(p,w)=\max_{\theta\in\mathcal C}
 \left\{\sum_i p_i(\theta_i+w_i)-\min_j(\theta_j+w_j)\right\}
 =\max_j\sum_i p_iH_{ji}(w).
\end{equation}
Thus the optimal allocation and its value $V(w)$ are the solution of a finite linear program. For any rational $p,\lambda\in\Delta_q$, the computable numbers
\[
 L=\min_i\sum_j\lambda_jH_{ji}(w),\qquad
 U=\max_j\sum_i p_iH_{ji}(w)
\]
satisfy $L\le V(w)\le U$. There exist optimal primal and dual vectors with $L=U$. On the original coverage event, every implemented allocation has expected catalogue regret at most $U$, simultaneously for all $p,w$, including choices made after examining the qualified intervals. The worst-case bound in \eqref{eq:allocation79} is attained by one of the cost vectors already constructed in Theorem~\ref{thm:choice78}.
\end{theorem}

This is a decision calculation from the purchased observations, not a new experiment or an assertion that uncertainty reduction changes the true controller costs. Pure choices remain available as vertices of $\Delta_q$. Consequently $V(w)\le\min_i\Gamma_i(w)$, with strict improvement possible when no pure controller is uniformly best. The comparator is the best single qualified controller at the true cost vector. The original Bellman-optimum guarantees and the empirical comparisons between fitted representations remain separate.

\paragraph{An exact acquisition-fee region.}
Let $A$ and $B$ denote the fine and coarse ReLU controllers. The qualified gain is $J_B-J_A\in[l,u]$, where $0<l<u$. All other controllers are dominated by $B$ when their charges are equal. Suppose only $A$ has the incremental per-deployment charge $f$. For $l<f<u$, the fine-controller share and sharp expected regret are
\begin{equation}\label{eq:two-allocation79}
 p_A^*(f)=\frac{u-f}{u-l},\qquad
 V(f)=\frac{(f-l)(u-f)}{u-l}.
\end{equation}
For $f\le l$, use $A$ alone; for $f\ge u$, use $B$ alone. At the fee midpoint, allocation halves the best pure worst-case regret: $(u-l)/4$ rather than $(u-l)/2$. This is an exact decision bound, not a newly measured expected-cost reduction. In the retained experiment these bounds are approximately $0.00024059$ and $0.00048118$, respectively; Table~\ref{tab:allocation79} gives outward endpoints from the original rational intervals.

Maintaining both controllers need not be free. If a nontrivial allocation incurs an additional normalized cost $\kappa$, not incurred by either pure choice, its advantage survives precisely when
\begin{equation}\label{eq:setup-threshold79}
 \kappa<\frac{\min\{f-l,u-f\}^{2}}{u-l}.
\end{equation}
The full installed-subset account and a deterministic finite-cohort implementation are proved in Appendix~\ref{app:allocation79}. They charge every selected controller and do not turn shared experimental qualification into an unmeasured cold-start saving. For $m$ deployments and monetary loss scale $v>0$, the two-controller setup allowance is $mv$ times the right side. This gives an economically explicit frontier conditional on the actual fee, loss scale and deployment count, without inventing their calibration.

\section{Conclusion}\label{sec:conclusion65}
Neural Bellman Operators connect attained continuation approximation to feasible policy decisions under the original economic law. The constructive readout, uniform qualification and direct controller-cost experiment form one chain; optional query prediction and witnessed reuse form distinct work mechanisms. Exact integration removes optimal-value descendants from evaluation of a qualified shallow continuation, while the original fitting, qualification and every failed reuse remain paid operations.

The matched continuation experiment identifies lower ReLU-controller costs at both acquisition resolutions. Its joint intervals also determine the constructed-policy ordering and conditional purchase regions. The deployment theorem extends those observations to sharp expected-regret allocations, additional installation costs and finite cohorts. It does not turn decision-theoretic regret reduction into a newly observed cost gain. Actual monetary conclusions require the stated loss scale, deployment count and resource charges. All historical adverse selector findings and broader economic applications remain available with their own assumptions and evidence.

\bibliographystyle{ecta-fullname}
\bibliography{references,references78}
\end{document}
